% 原书 p.1–p.3（PDF p.9–p.11）
% >>> p009
\chapter{辛流形（Symplectic Manifold）}

辛几何是几何与拓扑的一个分支，它研究辛流形及其子流形的几何性质、辛流形上的各种动力学行为，
以及不同辛流形之间的神秘联系。其核心目标之一，是发现由辛结构所施加的约束而产生的各种刚性现象。
本章将介绍这一学科的主要构件，它们为后续各章的进一步发展提供背景。

\section{辛结构（Symplectic Structure）}

\begin{keytakeaways}
\begin{itemize}
  \item 几何图像：辛结构是微分流形上的一个非退化闭二形式 $\omega$。它给出切空间到余切空间的典范同构
        $v\mapsto\omega_p(v,\cdot)$，从而把向量场与函数一一对应——这正是哈密顿力学的起点。
  \item 代数工具：斜对称双线性形式、外微分 $\dd$、楔幂 $\omega^k$ 与 de Rham 上同调类 $[\omega^k]$。
  \item 本节结论：辛流形必为偶数维、必可定向，且 $\omega^n$ 是无处零化的体积形式；闭辛流形上
        $[\omega^k]=[\omega]^k\neq0$（命题 1.1.1），故 $S^{2n}$（$n\ge2$）等流形不允许辛结构。
  \item 线性代数内核：命题 1.1.2 断言 $n-k$ 次形式与 $n+k$ 次形式被 $\omega^k$ 的楔积同构起来，
        这是纯线性代数的「强 Lefschetz 性质」，与 $M$ 的整体拓扑无关。
\end{itemize}
\end{keytakeaways}

\begin{wdefn}{1.1.1}
设 $M$ 为微分流形，$\omega\in\Omega^2(M)$ 为 $M$ 上的二形式。若 $\omega$ 满足以下条件，则称 $\omega$
为 $M$ 上的\textbf{辛结构}：
\begin{enumerate}[label=(\roman*)]
  \item $\omega$ 是\textbf{非退化的}：对任意 $p\in M$ 与非零的 $v\in T_pM$，由
        $v\mapsto\omega_p(v,\cdot)$ 给出的映射 $T_pM\to T_p^*M$ 是同构；
  \item $\omega$ 是\textbf{闭的}，即 $\dd\omega=0$。
\end{enumerate}
若如此，则称 $(M,\omega)$ 为一个\textbf{辛流形}（symplectic manifold）。
\end{wdefn}

由定义 1.1.1 可直接得到以下几条有用的观察：

\begin{enumerate}[label=(\alph*)]
  \item 定义 1.1.1 中的条件 (i) 表明：对任意非零的 $v\in T_pM$，总存在 $w\in T_pM$ 使得
        $\omega_p(v,w)\neq0$。至多差一个符号，数值 $\omega_p(v,w)$ 通常称为由向量 $v$ 与 $w$
        张成的区域的 $\omega$-\textbf{面积}（$\omega$-area）。
  \item 任何辛流形 $(M,\omega)$ 都是偶数维的。事实上，按定义 $\omega_p$ 是 $T_pM$ 上的斜对称双线性形式，
        而定义 1.1.1 中的条件 (i) 表明：在其标准形中不会出现落入核的部分。
  \item 任何辛流形 $(M,\omega)$ 都是可定向的（因此，例如莫比乌斯带就不允许任何辛结构）。事实上，
        设 $\dim M=2n$，则容易验证定义 1.1.1 中的条件 (i) 等价于
        $\omega^n=\omega\wedge\cdots\wedge\omega\in\Omega^{2n}(M)$
        是无处零化的形式。因此 $\omega^n$ 是 $M$ 上的体积形式，这蕴涵 $M$ 可定向。
\end{enumerate}

寻找流形允许辛结构的必要条件，是辛几何的主要问题之一。下面是一条有趣的必要条件。

% >>> p010
\begin{wprop}{1.1.1}
设 $(M,\omega)$ 为维数 $2n$ 的闭辛流形，则对任意 $1\le k\le n$，de Rham 上同调类
$[\omega^k]=[\omega]^k\neq0$ 位于 $H^{2k}(M;\R)$ 之中。
\end{wprop}

\begin{wproof}
反设对某个 $1\le k\le n$ 有 $[\omega^k]=0$，则存在某个 $(2k-1)$-形式 $\alpha\in\Omega^{2k-1}(M)$
使得 $\omega=\dd\alpha$。由于按定义 1.1.1 的条件 (ii)，$\omega$ 是闭的，故
$\omega^n=\dd(\alpha\wedge\omega^{n-k})$。于是由 Stokes 定理，
\[
  \int_M\omega^n=\int_M\dd(\alpha\wedge\omega^{n-k})
  =\int_{\partial M}\alpha\wedge\omega^{n-k}=0,
\]
其中最后一个等号来自 $M$ 是闭的（即无边界）这一假设。这与上面观察 (c) 中「$\omega^n$ 是 $M$ 上的
体积形式」相矛盾。于是得到所需的结论。
\end{wproof}
% [Tutor Note: 原文疑误——此处「存在 (2k−1)-形式 α 使得 ω = dα」应为「使得 ω^k = dα」：由 [ω^k] = 0
% 只能推出 ω^k 恰当；而由条件 (i) 非退化，ω 本身在闭流形上绝不可能是恰当的（否则 ω^n 的积分会为 0，
% 与体积形式矛盾）。后续 ω^n = d(α ∧ ω^{n−k}) 也只有在 α 满足 dα = ω^k 时才成立。译文保留原文写法，
% 详见 review_ch01.md。另注：原文此句句首的逗号后接大写「Then」，属印刷失误。]

\begin{wrem}{1.1.1}
命题 1.1.1 说明：辛流形的偶数维 de Rham 上同调群在很强的意义下全都非零——它们含有由同一个二形式
的楔幂给出的上同调类。
\end{wrem}

本命题立即排除掉许多闭流形（不允许辛结构），即便它们是偶数维且可定向的。

\begin{wexam}{1.1.1}
任何偶数维球面 $S^{2n}$（$n\ge2$）都没有辛结构，因为当 $n\ge2$ 时 $H^2(S^{2n};\R)=0$。\hfill$//$
\footnote{在全书范围内，每个例子末尾的记号 // 表示该例结束。}
\end{wexam}

\begin{wexer}{1.1.1}
证明：两个球面的积 $S^p\times S^q$ 允许辛结构，当且仅当 $p=q=1$ 或 $p=q=2$（提示：利用命题 1.1.1
与注 1.1.1）。
\end{wexer}

\begin{wque}{1.1.1}
若 $M$ 是闭的、$4$ 维、可定向的光滑流形且 $H^2(M;\R)\neq0$，那么 $M$ 是否总允许辛结构？
\end{wque}

不过，问题 1.1.1 的答案是否定的。一个例子是连通和 $\CP^2\#\CP^2\#\CP^2$，它由 C. Taubes [214]
借助 Seiberg--Witten 理论（超出本书范围）所发现。
% [Tutor Note: 注意此处的例子是三个 CP² 的连通和（定向相同），而不是 CP² 与 \bar{CP}² 的连通和：
% 后者是 CP² 的爆破，仍然是辛流形；而 CP²#CP²#CP² 的交形式为奇型且 b_2^+=3，Taubes 用
% Seiberg–Witten 不变量证明它不能承载辛结构。这是「H²(M;R) ≠ 0 且可定向」并不足以保证辛结构的
% 标准反例。]

注意命题 1.1.1 只对闭（辛）流形有效。作为本节的结束，我们再提一个局部有效的必要条件（即：不需要
定义 1.1.1 的条件 (ii) 即可确认）。它模仿了复几何中的强 Lefschetz 性质（[111] 中的定义 3.8），
只不过是在形式层面上。

\begin{wprop}{1.1.2}（[38] 中的引理 1.5）
设 $(M,\omega)$ 为维数 $2n$ 的辛流形。则对任意 $p\in M$ 与任意 $0\le k\le n$，诱导映射
$\wedge\omega_p^k:\wedge^{n-k}T_p^*M\to\wedge^{n+k}T_p^*M$ 是同构。
\end{wprop}

% >>> p011
\begin{wproof}
这是一个纯代数性质，故为简化记号记 $V=T_p^*M$。由维数计算，
\[
  \dim\wedge^{n-k}V=\binom{2n}{n-k}=\binom{2n}{n+k}=\dim\wedge^{n+k}V .
\]
因此只需证明诱导映射 $\wedge\omega_p^k:\wedge^{n-k}V\to\wedge^{n+k}V$ 是单射。当 $k=n$ 时，
由定义 1.1.1 的条件 (i)（或上面的观察 (c)），诱导映射
$\wedge\omega_p^n:\R(=\wedge^0V)\to\R(=\wedge^{2n}V)$ 是单射，因为它只是一个非零数的乘法。
又当 $k=0$ 时，诱导映射就是恒等映射，其单性显然。

现在，设对某个 $k\in\{1,\cdots,n\}$，$\wedge\omega_p^k:\wedge^{n-k}V\to\wedge^{n+k}V$ 的单性成立。
则若对某个 $\alpha\in\wedge^{n-k+1}V$ 有 $\alpha\wedge\omega_p^{k-1}=0$，就有
$\alpha\wedge\omega_p^k=(\alpha\wedge\omega_p^{k-1})\wedge\omega_p=0$。对任意向量
$Y\in T_pM=V^*$，我们有
\[
  0=\iota_Y(\alpha\wedge\omega_p^k)
  =\iota_Y\alpha\wedge\omega_p^k+k\cdot\alpha\wedge\omega_p^{k-1}\wedge\iota_Y\omega .
\]
这说明 $\iota_Y\alpha\wedge\omega_p^k=0$。由假设 $\wedge\omega_p^k:\wedge^{n-k}V\to\wedge^{n+k}V$
是单射，故 $\iota_Y\alpha=0$。由于这对任意 $Y$ 都成立，我们得到 $\alpha=0$。于是由归纳论证
（从 $k=n$ 倒推至 $k=1$）即得所需结论。
\end{wproof}
% [Tutor Note: 此处 $V=T_p^*M$，故 $V^*=T_pM$，$\iota_Y$ 是沿切向量 $Y$ 的内乘缩并，它作用在
% $\wedge^\bullet V=\wedge^\bullet T_p^*M$ 上并把次数降低 1。缩并公式中的系数 $k$ 来自 Leibniz 法则：
% 对 $k$ 个因子的 $\omega$ 逐个缩并，共得到 $k$ 个相同的项。易错点：$Y$ 是切向量（而非余切向量），
% 且 $\iota_Y\omega=\omega(Y,\cdot)$ 按本书约定（见 (1.15)）与 $\omega(\cdot,Y)$ 差一个符号。]

\begin{wrem}{1.1.2}
当然，也可以在上同调类层面考虑原来的强 Lefschetz 性质，即由 $a=[\alpha]\mapsto[\alpha\wedge\omega^k]$
给出的映射 $H^{n-k}(M;\R)\to H^{n+k}(M;\R)$。遗憾的是，命题 1.1.2 的论证无法以同样的方式启动：
因为先验地我们并不知道 $H^{n-k}(M;\R)$ 与 $H^{n+k}(M;\R)$ 之间的任何关系。有趣的是，这一观察在历史上
引出了区分辛流形与 Kähler 流形的尝试。
\end{wrem}
